\newpage
\begin{center}
\large\textbf{Минобрнауки России}

\large\textbf{Юго-Западный государственный университет}
\vskip 1em
\normalsize{Кафедра программной инженерии}
\vskip 1em
\ifВКР{
        \begin{flushright}
        \begin{tabular}{p{.4\textwidth}}
        \centrow УТВЕРЖДАЮ: \\
        \centrow Заведующий кафедрой \\
        \hrulefill \\
        \setarstrut{\footnotesize}
        \centrow\footnotesize{(подпись, инициалы, фамилия)}\\
        \restorearstrut
        «\underline{\hspace{1cm}}»
        \underline{\hspace{3cm}}
        20\underline{\hspace{1cm}} г.\\
        \end{tabular}
        \end{flushright}
        }\fi
\end{center}
\vspace{1em}
  \begin{center}
  \large
\ifВКР{
ЗАДАНИЕ НА ВЫПУСКНУЮ КВАЛИФИКАЦИОННУЮ РАБОТУ
  ПО ПРОГРАММЕ БАКАЛАВРИАТА}
  \else
ЗАДАНИЕ НА КУРСОВУЮ РАБОТУ (ПРОЕКТ)
\fi
\normalsize
  \end{center}
\vspace{1em}
{\parindent0pt
  Студента \АвторРод, шифр\ \Шифр, группа \Группа
  
1. Тема «\Тема\ \ТемаВтораяСтрока\ \ТемаТретьяСтрока»
\ifВКР{
утверждена приказом ректора ЮЗГУ от \ДатаПриказа\ № \НомерПриказа
}\fi.

2. Срок предоставления работы к защите \СрокПредоставления

3. Исходные данные для создания программной системы:

3.1. Перечень решаемых задач:}

\renewcommand\labelenumi{\theenumi)}

\begin{enumerate}
\item проанализировать предметную область навигации внутри зданий и определить требования к системе;
\item разработать архитектуру программного комплекса для построения маршрутов внутри здания с учетом этажности, переходов и ограничений на перемещение;
\item реализовать средства работы с навигационным графом, планами этажей, QR-навигацией и расписанием;
\item создать интерфейс пользователя и административные инструменты для редактирования карты, маршрутов и справочных данных;
\item предусмотреть поддержку режима чрезвычайной ситуации с построением эвакуационных маршрутов;
\item проверить корректность работы основных сценариев использования и подготовить материалы по результатам разработки.
\end{enumerate}

{\parindent0pt
  3.2. Входные данные и требуемые результаты для программы:}

\begin{enumerate}
\item Входными данными для программной системы являются: планы этажей в формате SVG, сведения о зданиях, этажах и навигационных узлах, навигационный граф, данные пользователей и их ролей, расписание занятий, параметры маршрутизации, а также состояние режима чрезвычайной ситуации.
\item Выходными данными для программной системы являются: построенный маршрут между выбранными точками, пошаговые инструкции перемещения, визуальное отображение маршрута на плане этажа, QR-данные для навигационных узлов, сведения о текущем расписании и эвакуационные маршруты при активации режима чрезвычайной ситуации.
\end{enumerate}

{\parindent0pt

  4. Содержание работы (по разделам):
  
  4.1. Введение.
  
4.2. Анализ предметной области: назначение систем построения эвакуационных
маршрутов в реальном времени для замкнутых пространств, пользовательские
роли, основные сценарии определения местоположения, особенности работы с
картой здания, QR-навигацией, заблокированными участками и режимом
чрезвычайной ситуации.
  
4.3. Техническое задание: основание для разработки, назначение системы
построения эвакуационных маршрутов, требования к авторизации, построению
маршрутов, отображению карты, QR-навигации, административным функциям,
режиму чрезвычайной ситуации и оформлению программной документации.

4.4. Технический проект: общая характеристика системы построения
эвакуационных маршрутов, обоснование выбора {.NET MAUI}, XAML и
SkiaSharp, описание архитектуры MVVM, моделей навигационного графа,
сервисного слоя, пользовательского интерфейса, алгоритма построения маршрута
и работы с локальными JSON-данными.

4.5. Рабочий проект: спецификация основных компонентов и классов приложения,
описание экранов пользователя и администратора, проверка сценариев
авторизации, выбора аудитории, построения маршрута, сканирования QR-кода,
редактирования карты и активации режима чрезвычайной ситуации.

4.6. Заключение.

4.7. Список использованных источников.

5. Перечень графического материала:

\списокПлакатов

\vskip 2em
\begin{tabular}{p{6.8cm}C{3.8cm}C{4.8cm}}
Руководитель \ifВКР{ВКР}\else работы (проекта) \fi & \lhrulefill{\fill} & \fillcenter\Руководитель\\
\setarstrut{\footnotesize}
& \footnotesize{(подпись, дата)} & \footnotesize{(инициалы, фамилия)}\\
\restorearstrut
Задание принял к исполнению & \lhrulefill{\fill} & \fillcenter\Автор\\
\setarstrut{\footnotesize}
& \footnotesize{(подпись, дата)} & \footnotesize{(инициалы, фамилия)}\\
\restorearstrut
\end{tabular}
}

\renewcommand\labelenumi{\theenumi.}
